\NSPage{103}%
For two curl-generated waves \NSi{NS-F-P103-001}{w,w'} with the same label and exponents
\NSi{NS-F-P103-002}{\alpha,\alpha'}, the nonzero harmonics of
\NSi{NS-F-P103-003}{(w\mathbin{\cdot}\nabla_*)w'+(w'\mathbin{\cdot}\nabla_*)w} belong to
\NSi{NS-F-P103-004}{\mathsf W_{\alpha+\alpha'-\kappa_s}}; the zero harmonic belongs to the mean
class \NSi{NS-F-P103-005}{\mathsf M_{\alpha+\alpha'-\kappa_s}}. Distinct labels have zero products
on their closed supports.
\end{lemma}

\begin{proof}
Mean advection of a wave phase costs \NSi{NS-F-P103-006}{k=O(\varepsilon^{-1/2})}. Its radial
transport of the amplitude has a radial mean factor of order \NSi{NS-F-P103-007}{\mu+1} and one
\NSi{NS-F-P103-008}{D_r}; axial transport gains one through \NSi{NS-F-P103-009}{D_z}. Wave
transport of a mean costs at most \NSi{NS-F-P103-010}{\kappa_s} radially, and the cylindrical
connections contain no derivatives. These terms are all bounded by the asserted mean--wave
exponent.

For two harmonics of one label, phase differentiation in the transport of
\NSi{NS-F-P103-011}{a'\exp(ikm'\Phi)} by \NSi{NS-F-P103-012}{a\exp(ikm\Phi)} gives
\NSi{NS-F-P103-013}{ikm'(a\mathbin{\cdot}n_\Phi)a'}. Equation~\eqref{eq:7.39} cancels the
apparent half-power loss. Transport terms in which the derivative acts on the amplitude cost at
most \NSi{NS-F-P103-014}{\kappa_s}, and the connections cost none. Products of the weights obey
\NSi{NS-F-P103-015}{P_v^2\le P_v} and have radial weight \NSi{NS-F-P103-016}{\zeta}; this is also
stronger than the required \NSi{NS-F-P103-017}{\sqrt\zeta} weight for a nonzero output. The support
separation in Lemma~\ref{lem:6.1} proves the last assertion. Leibniz' rule proves each statement for
amplitude derivatives through every fixed order.
\end{proof}

\subsection{Sources and errors retained during iteration.}\label{sec:9.2}
To use the estimates in Proposition~\ref{prop:9.1} and Lemma~\ref{lem:9.2} at successive stages,
each new source must satisfy the support and derivative hypotheses of the pulse inverse in
Proposition~\ref{prop:7.2}. Mean operations can enlarge auxiliary support, so we must check these
hypotheses after a complete cycle. We keep the errors already flat at the singularity separate from
the sources to be corrected, as in \eqref{eq:9.3}, while using the full velocity fields in every
nonlinear product. The supported sources must remain inside the fixed enlarged rectangle;
Proposition~\ref{prop:9.3} establishes these properties.

We use the source supports \NSi{NS-F-P103-018}{\Omega_\gamma} and enlarged domains
\NSi{NS-F-P103-019}{\mathcal E_\gamma^+} of \eqref{eq:6.28}. The flat term below collects the base
error \NSi{NS-F-P103-020}{\mathcal E_B} from \eqref{eq:5.41}, pulse-cutoff tails, and uncompensated
cutoff remainders. The exact mean balances can contain such flat terms; their definition always uses
the complete current velocity. Angular means may occupy the full auxiliary torus; the rectangle and
envelope conditions below concern the nonzero harmonics.

\begin{proposition}\label{prop:9.3}
Fix \NSi{NS-F-P103-021}{j\in\mathbb N_0}. Let
\NSi{NS-F-P103-022}{(u^{[j]},p^{[j]})} be obtained from the fixed slow base and primary waves by a
finite sequence of the pulse inverse and curl construction, the signed amplitude map for
auxiliary-independent stresses, and the mean maps of Section~\ref{sec:8}. Suppose that pressure is
reconstructed by \eqref{eq:8.12} after each update and all products use the complete velocity fields.
There is a physical residual decomposition
\NSFormula{NS-F-P103-023}{display}{9.3}
\begin{equation}
\mathscr R(u^{[j]},p^{[j]})=\mathcal G^{[j]}+\mathcal F^{[j]}.
\tag{9.3}\label{eq:9.3}
\end{equation}
In a common chart, put
\NSi{NS-F-P103-024}{\mathcal G_*^{[j]}=Q^{2A+1/2}\mathcal G^{[j]}} and
\NSi{NS-F-P103-025}{\mathcal F_*^{[j]}=Q^{2A+1/2}\mathcal F^{[j]}}. Then:
\begin{enumerate}[label=(\roman*)]
\item \emph{Supported sources.} On the coarsest common auxiliary torus, the nonzero angular
harmonics have the form
\NSFormula{NS-F-P103-026}{display}{9.4}
\begin{equation}
\begin{gathered}
\mathcal G_*^{[j]}-\langle\mathcal G_*^{[j]}\rangle_\theta
 =\sum_{\gamma\in\Gamma}\sum_{m\in\mathscr H_j}
 f_{\gamma,m}^{[j]}e^{ik_\gamma m\Phi_\gamma},\\
\mathscr H_j\subset\mathbb Z\setminus\{0\},\qquad |\mathscr H_j|<\infty.
\end{gathered}
\tag{9.4}\label{eq:9.4}
\end{equation}
The sum is locally finite, and \NSi{NS-F-P103-027}{\mathscr H_j} depends on the finite construction
but not on the band. The coefficients satisfy
\NSFormula{NS-F-P103-028}{display}{none}
\[
\operatorname{supp}(f_{\gamma,m}^{[j]}|_{\mathcal E_\gamma})\subset\Omega_\gamma,
\qquad
f_{\gamma,m}^{[j]}\in C^\infty(\mathcal E_\gamma^+;\mathbb C^3),
\]
with smooth zero extension across the fixed slow and transverse supports and the enlarged local
torus rectangle. Let \NSi{NS-F-P103-029}{\alpha,\mu} be any lower bounds for the wave and mean
exponents furnished by
\NSPage{104}%
applying Propositions~\ref{prop:9.1}, \ref{prop:7.2}, \ref{prop:6.6} and \ref{prop:7.6},
Section~\ref{sec:8},
and Lemma~\ref{lem:9.2} along this finite sequence. Then
\NSFormula{NS-F-P104-001}{display}{none}
\[
f_{\gamma,m}^{[j]}\in \mathsf W_\alpha,
\qquad
\langle\mathcal G_*^{[j]}\rangle_\theta\in \mathsf M_\mu.
\]
The mean coefficients have smooth zero extension outside the active shell. The next pulse inverse is
applied precisely to the coefficients \NSi{NS-F-P104-002}{f_{\gamma,m}^{[j]}} in \eqref{eq:9.4}.

\item \emph{Flat remainder.} For every fixed \NSi{NS-F-P104-003}{j,m,N},
\NSFormula{NS-F-P104-004}{display}{9.5}
\begin{equation}
|\mathcal F^{[j]}|_m\le C_{j,m,N}q^N.
\tag{9.5}\label{eq:9.5}
\end{equation}
This physical bound is uniform over bands, labels, and points.

\item \emph{Exact mean equations.} Let
\NSi{NS-F-P104-005}{(\beta^{[j]},v^{[j]},\gamma^{[j]},w^{[j]})} be the complete current velocity
correction as in \eqref{eq:8.1}, and set
\NSi{NS-F-P104-006}{\mathcal E_{B,*}=Q^{2A+1/2}\mathcal E_B}. Then
\NSFormula{NS-F-P104-007}{display}{9.6}
\begin{equation}
\begin{gathered}
W_{ab}^{[j]}=\langle w_a^{[j]}w_b^{[j]}\rangle_\theta,
\qquad
\mathcal P^{[j]}=\int\langle g_r^{[j]}\rangle_Y\,dR,
\qquad
p_m^{[j]}=\mathsf T_0(g_r^{[j]}-\rho\mathcal P^{[j]}),\\
\langle\mathcal G_*^{[j]}+\mathcal F_*^{[j]}-\mathcal E_{B,*}\rangle_\theta
=\bigl(D_r p_m^{[j]}-g_r^{[j]},E_\theta^{[j]},E_z^{[j]}\bigr),
\end{gathered}
\tag{9.6}\label{eq:9.6}
\end{equation}
where \NSi{NS-F-P104-008}{g_r^{[j]},E_\theta^{[j]},E_z^{[j]}} are given by \eqref{eq:8.3} on
this tuple.
\end{enumerate}
\end{proposition}

\begin{proof}
We first prove (i), then track the additive errors in (ii) while preserving the exact mean equations
in (iii).

\noindent\textbf{Part (i): support and coefficient bounds.} The primary pulses have the claimed
supports and weighted bounds at every derivative order. Differentiation preserves closed support and
the amplitude class bounds. An explicit \NSi{NS-F-P104-009}{D_r} costs
\NSi{NS-F-P104-010}{\kappa_s} each time it occurs; taking further amplitude derivatives of that
estimate enlarges its constants and polynomial degrees, without repeating that exponent loss.

For products, a mean times a wave inherits the wave's slow, transverse, and local torus rectangle
support, and its \NSi{NS-F-P104-011}{P_v} factor. Two waves of one label add their integer phase
multipliers. Their zero harmonic is exactly their angular mean, since the angular frequency
\NSi{NS-F-P104-012}{kp} is a nonzero integer; every other output retains the wave weight by
\NSi{NS-F-P104-013}{P_v^2\le P_v}. Different labels do not interact. Curls, pressure coefficients,
and particular pulse propagation preserve the phase integer. Mean operations cannot introduce a
nonzero angular harmonic. Any subsequent nonzero harmonic product contains a wave factor and
therefore inherits that factor's support and envelope.

The radial integration, pressure reconstruction, mean-vector-potential construction, and
auxiliary-time inversion in Section~\ref{sec:8} preserve all mean-amplitude bounds. For the stress correction
entering Proposition~\ref{prop:7.6}, subtract a supported interior profile with the same weighted
radial moment as the source. The resulting weighted radial integral can be taken forward from the
left support boundary or backward from the right support boundary, because the corrected source has
zero total moment. Near either boundary, let \NSi{NS-F-P104-014}{s} be the corresponding logarithmic
distance in \eqref{eq:6.23}. The weight \NSi{NS-F-P104-015}{\zeta} is
\NSi{NS-F-P104-016}{e^{-a/s^2}} times a smooth positive factor bounded above and below there, with
\NSi{NS-F-P104-017}{a>0} fixed. For each fixed \NSi{NS-F-P104-018}{M\ge0} and integer
\NSi{NS-F-P104-019}{k\ge0}, the following bounds hold for sufficiently small positive
\NSi{NS-F-P104-020}{\delta,s}:
\NSFormula{NS-F-P104-021}{display}{none}
\[
\int_0^\delta s^{-M}e^{-a/s^2}\,ds
 \le C_{M,a}\delta^{3-M}e^{-a/\delta^2},
\qquad
\left|\frac{d^k}{ds^k}e^{-a/s^2}\right|
 \le C_{k,a}s^{-3k}e^{-a/s^2}.
\]
\NSPage{105}%
The first bound follows by the substitution \NSi{NS-F-P105-001}{v=a/s^2}; the second follows by
differentiation. Radial weights and the radial Jacobian are bounded on the normalized shell.
Differentiating the radial integral therefore retains
\NSi{NS-F-P105-002}{\zeta\delta^{-M'}}, with \NSi{NS-F-P105-003}{M'} allowed to depend on the
derivative; in the interior \NSi{NS-F-P105-004}{\zeta} is bounded below. Thus a source of order
\NSi{NS-F-P105-005}{\alpha} gives a stress correction
\NSi{NS-F-P105-006}{\Sigma\in \mathsf M_\alpha}.

For the amplitude division, use the same normalization as \eqref{eq:7.31}:
\NSi{NS-F-P105-007}{d_\Sigma=\mathsf H^{-1}(\Sigma/\varepsilon)} and
\NSi{NS-F-P105-008}{\delta a_\sigma=(d_\Sigma)_\sigma/(2a_\sigma)}. The inverse-matrix bounds and
\eqref{eq:7.25} give, for every fixed coefficient multi-index \NSi{NS-F-P105-009}{I},
\NSFormula{NS-F-P105-010}{display}{none}
\[
\begin{aligned}
|D^I(d_\Sigma)_\sigma|
 &\le C_I\varepsilon^{\alpha-1}S_*^{b_I}\zeta\delta^{-M_I},\\
|D^I(a_\sigma^{-1})|
 &\le C_I S_*^{b_I}\zeta^{-1/2}\delta^{-M_I},\\
|D^I\delta a_\sigma|
 &\le C_I\varepsilon^{\alpha-1}S_*^{b'_I}\sqrt\zeta\,\delta^{-M'_I},\\
|D^I(\sqrt\varepsilon\,\delta a_\sigma b_\sigma)|
 &\le C_I\varepsilon^{\alpha-1/2}S_*^{b''_I}\sqrt\zeta\,\delta^{-M''_I}P_v.
\end{aligned}
\]
This recovers \eqref{eq:7.33}. The last estimate uses the fundamental pulse bounds; after the slow
cutoff it is the \NSi{NS-F-P105-011}{\mathsf W_{\alpha-1/2}} estimate of
Proposition~\ref{prop:7.6}. The remaining square-root edge weight gives smooth extension by zero.

The hypotheses of Proposition~\ref{prop:7.2} now hold for each new source. Its prescribed path holds
the slow and transverse variables fixed; zero data along an entire such path give zero solution. The
final cutoff restores the fixed enlarged local torus rectangle support. Smooth source extensions give
smooth extensions of the solution. This proves the support-containment hypothesis of the inverse;
divisibility by the original cutoff is unnecessary.

The local band set remains fixed: radial paths hold \NSi{NS-F-P105-012}{(z,t)} fixed, and
\NSi{NS-F-P105-013}{q} depends only on \NSi{NS-F-P105-014}{(z,t)}; temporal inversion and pulse
propagation hold all slow variables fixed. Hence the common-torus construction of
Lemma~\ref{lem:6.2} applies at every step of the finite construction. A radial mean can aggregate
polynomially many labels in \NSi{NS-F-P105-015}{S_*}, whereas only boundedly many wave labels
overlap at a point. Radial aggregation changes only the polynomial factors in
\NSi{NS-F-P105-016}{S_*}, while preserving the common chart and the
\NSi{NS-F-P105-017}{\varepsilon} exponent.

\noindent\textbf{Parts (ii) and (iii): flat errors and exact means.} For an actual increment
\NSi{NS-F-P105-018}{(w,\pi)}, the residual changes by
\NSFormula{NS-F-P105-019}{display}{none}
\[
\begin{aligned}
\mathscr R(u+w,p+\pi)={}&\mathscr R(u,p)+\partial_t w-\Delta w+\nabla\pi\\
&+(u\mathbin{\cdot}\nabla)w+(w\mathbin{\cdot}\nabla)u+(w\mathbin{\cdot}\nabla)w.
\end{aligned}
\]
The old \NSi{NS-F-P105-020}{\mathcal F} enters additively, while all displayed products use actual
fields, including their cutoff curls and remainders from the compactly supported modified radial
integral. Add newly created flat terms to \NSi{NS-F-P105-021}{\mathcal F}. At a fixed stage there
are finitely many kinds of such terms. The Gaussian bound in Proposition~\ref{prop:9.1} is uniform
over labels. The estimate in Lemma~\ref{lem:8.2} bounds each cutoff remainder by any prescribed
power of \NSi{NS-F-P105-022}{q}, using finitely many additional source derivatives. Summing
polynomially many labels preserves these bounds. Keep Gaussian nonzero-harmonic cutoff tails in the
additive residual and exclude them from subsequent forward pulse inputs. The mean residuals
\NSi{NS-F-P105-023}{E_\theta,E_z} are the exact conservative balances and may contain flat cutoff
remainders, which a fixed mean inverse may cancel. Remove any canceled contribution from
\NSi{NS-F-P105-024}{\mathcal F}. The cutoff-remainder bounds place the remainder in
\NSi{NS-F-P105-025}{\mathsf M_\alpha} for every \NSi{NS-F-P105-026}{\alpha}, so this update
preserves the stated classes and the exact residual decomposition. This proves \eqref{eq:9.5} at each
finite stage, hence (ii). Applying Proposition~\ref{prop:8.1} to the complete current tuple, and using
\eqref{eq:8.12}, gives \eqref{eq:9.6} and proves (iii).
\end{proof}

The realization step will use a comparison with one fixed finite stage; it will not sum these flat
residuals directly.

\NSPage{106}%
\subsection{Initialization and a full correction cycle.}\label{sec:9.3}
We now put the wave and mean corrections together. The primary pulses, followed by the initial mean
corrections, will give the first set of residual bounds in Proposition~\ref{prop:9.5}. We then show in
Proposition~\ref{prop:9.6} that a complete cycle improves all these bounds while preserving the
support and moment conditions. The induction statement below combines the residual orders in
\eqref{eq:9.8} with the bounds \eqref{eq:9.9} on the total velocity correction that control its
interactions with later increments.

We keep the base \NSi{NS-F-P106-001}{(b,V,G)}, the primary amplitudes, and all inverse operators
fixed throughout the construction. Let
\NSi{NS-F-P106-002}{w_0^{\mathrm{tan}}=W_0^{\mathrm{as}}} be the normalized representative of the
assembled primary transverse field in \eqref{eq:7.35}. For each finite construction write its
normalized velocity and pressure as
\NSFormula{NS-F-P106-003}{display}{9.7}
\begin{equation}
u_*=(b+\beta,V+v,G+\gamma)+w,
\qquad
p_*=p_{B,*}+p_m+p_w,
\qquad
\langle w\rangle_\theta=\langle p_w\rangle_\theta=0.
\tag{9.7}\label{eq:9.7}
\end{equation}
Here \NSi{NS-F-P106-004}{p_{B,*}} is the base pressure,
\NSi{NS-F-P106-005}{p_w} consists of the pressure harmonics supplied by the amplitude equations,
and \NSi{NS-F-P106-006}{p_m} is reconstructed from the current
\NSi{NS-F-P106-007}{g_r} by \eqref{eq:8.12}. The quantities
\NSi{NS-F-P106-008}{E_\theta,E_z,\mathcal P,\mathcal J_\theta,\mathcal J_z} are then determined by
\eqref{eq:8.3}, \eqref{eq:8.12}, and \eqref{eq:8.15}.

\begin{definition}[Finite correction state]\label{def:9.4}
For \NSi{NS-F-P106-009}{j\in\mathbb N\cup\{0\}}, a state at stage
\NSi{NS-F-P106-010}{j} consists of real fields of the form \eqref{eq:9.7} with the following
properties. The wave \NSi{NS-F-P106-011}{w} is a sum of curls of supported wave potentials with the
fixed label phases, and the mean correction \NSi{NS-F-P106-012}{(\beta,v,\gamma)} is the sum of a
curl of an azimuthal potential and a direct azimuthal field. In particular, both velocity corrections
are exactly divergence-free. The source supports, finite harmonic sets, smooth zero extensions, and
physical residual decomposition are as in Proposition~\ref{prop:9.3}.

The normalized nonzero harmonics \NSi{NS-F-P106-013}{\mathcal G_{\mathrm{wave}}^{[j]}} of
\NSi{NS-F-P106-014}{\mathcal G_*^{[j]}}, the exact tangential mean residuals, and the three integral
defects obey
\NSFormula{NS-F-P106-015}{display}{9.8}
\begin{equation}
\begin{gathered}
\mathcal G_{\mathrm{wave}}^{[j]}\in \mathsf W_{B_j},
\qquad E_\theta,E_z\in \mathsf M_{C_j^*},
\qquad (\mathcal P,\mathcal J_\theta,\mathcal J_z)\in \mathsf S_{C_j^*},\\
B_j=\frac12+\sigma_j,
\qquad C_j^*=1+\sigma_j,
\qquad \sigma_j=\frac15+\frac{j}{10},
\end{gathered}
\tag{9.8}\label{eq:9.8}
\end{equation}
with the following bounds for the total velocity and pressure corrections
\NSFormula{NS-F-P106-016}{display}{9.9}
\begin{equation}
w\in \mathsf W_{1/2},
\qquad w-w_0^{\mathrm{tan}}\in \mathsf W_{0.68},
\qquad v,\gamma,p_m\in \mathsf M_{0.9},
\qquad \beta\in \mathsf M_{1.9}.
\tag{9.9}\label{eq:9.9}
\end{equation}
Wave assertions refer to each grouped nonzero harmonic per label. The two normalized
angular-momentum and axial-flux moments are maintained exactly:
\NSFormula{NS-F-P106-017}{display}{9.10}
\begin{equation}
\int_0^\infty R^2\langle v\rangle_Y\,dR=0,
\qquad
\int_0^\infty R\langle\gamma\rangle_Y\,dR=0.
\tag{9.10}\label{eq:9.10}
\end{equation}
The pressure reconstruction is part of the state.

Consequently, the radial residual of every such state is
\NSi{NS-F-P106-018}{-\rho\mathcal P} plus the flat cutoff remainder of
Proposition~\ref{prop:8.3}.
\end{definition}

\begin{proposition}\label{prop:9.5}
Starting with the fixed base and the curls of the primary potentials, perform the temporal mean
update \eqref{eq:8.20} and the five-equation correction \eqref{eq:8.25}, reconstructing pressure by
\eqref{eq:8.12} before and after each update. The resulting fields form a state at stage
\NSi{NS-F-P106-019}{0} in the sense of Definition~\ref{def:9.4}; in particular, their residual orders
are \NSi{NS-F-P106-020}{B_0=0.7} and \NSi{NS-F-P106-021}{C_0^*=1.2}.
\end{proposition}

\begin{proof}
We first estimate the primary waves, then remove the torus-dependent means and the three moment
defects. The initialization uses the higher decay order of the auxiliary-averaged primary residual.
The primary transverse amplitudes have exponent \NSi{NS-F-P106-022}{1/2}. Their supported linear
residual components have exponent \NSi{NS-F-P106-023}{1-3\kappa_s}, and their nonlinear wave
residuals have
\NSPage{107}%
exponent \NSi{NS-F-P107-001}{1-\kappa_s}. The exact mean balances \eqref{eq:8.3} give
\NSi{NS-F-P107-002}{g_r,p_m,E_\theta,E_z\in \mathsf M_{1-\kappa_s}} and defects in
\NSi{NS-F-P107-003}{\mathsf S_{1-\kappa_s}} before any mean update.

For the auxiliary means, the primary covariance of Proposition~\ref{prop:7.5}, assembled by
\eqref{eq:7.35}, cancels the order-zero auxiliary radial--tangential stress tensor exactly, including
derivatives of the slow partition because the squared partition functions sum to one. To display this
cancellation, write \NSi{NS-F-P107-004}{\Sigma_a^{(0)}} for that leading tensor and
\NSi{NS-F-P107-005}{\Sigma_a} for the full tensor in \eqref{eq:8.3}. Before the mean velocities are
added, those equations read
\NSFormula{NS-F-P107-006}{display}{9.11}
\begin{equation}
\begin{aligned}
\langle E_\theta\rangle_Y
={}&(\partial_R+2/R)(\langle W_{r\theta}\rangle_Y-\Sigma_\theta^{(0)})
 -(\partial_R+2/R)(\Sigma_\theta-\Sigma_\theta^{(0)})
 +\varepsilon\partial_Z\langle W_{z\theta}\rangle_Y,\\
\langle E_z\rangle_Y
={}&(\partial_R+1/R)(\langle W_{rz}\rangle_Y-\Sigma_z^{(0)})
 -(\partial_R+1/R)(\Sigma_z-\Sigma_z^{(0)})
 +\varepsilon\partial_Z\langle W_{zz}+p_m\rangle_Y.
\end{aligned}
\tag{9.11}\label{eq:9.11}
\end{equation}
The first differences contain at least one curl correction. The tensor cross interaction between the
curl correction and a primary transverse amplitude belongs to
\NSi{NS-F-P107-007}{\mathsf M_{3/2-\kappa_s}} before applying the divergence. Axial fluxes, the
axial pressure term, and the higher-order auxiliary stress tensor have residual exponent at least
\NSi{NS-F-P107-008}{2-\kappa_s}. Thus
\NSi{NS-F-P107-009}{\langle E_\theta\rangle_Y,\langle E_z\rangle_Y\in \mathsf M_{1.49}}.

Set \NSi{NS-F-P107-010}{H_0=1-\kappa_s} and apply the temporal update \eqref{eq:8.20}. The
tangential increment is \NSi{NS-F-P107-011}{\mathsf M_{H_0}} and its induced radial velocity
component is \NSi{NS-F-P107-012}{\mathsf M_{H_0+1}}. The pressure change has exponent
\NSi{NS-F-P107-013}{H_0}, since \NSi{NS-F-P107-014}{2V\Delta v/R} appears in
\NSi{NS-F-P107-015}{\Delta g_r}. Its contribution to each tangential residual includes an axial
derivative and therefore has the corresponding additional gain. Beyond the cancelled fast-time
term, slow-time differentiation supplies one additional power of \NSi{NS-F-P107-016}{\varepsilon},
radial linear fluxes contain \NSi{NS-F-P107-017}{b=O(\varepsilon)} in every fixed-order amplitude
derivative on the shell or the induced radial velocity increment, axial fluxes gain one, and viscosity
gains \NSi{NS-F-P107-018}{1-2\kappa_s}. Radial quadratic mean fluxes also contain a radial mean.
All these residual changes have exponent at least
\NSi{NS-F-P107-019}{H_0+1-2\kappa_s}. Mean interaction with
\NSi{NS-F-P107-020}{w\in \mathsf W_{1/2}} has wave exponent at least
\NSi{NS-F-P107-021}{H_0}. Defects remain in \NSi{NS-F-P107-022}{\mathsf S_{H_0}} after pressure
is recomputed.

Now use the five-equation correction map \eqref{eq:8.25} at order
\NSi{NS-F-P107-023}{H_0} and recompute pressure. Its slow increments have no fast-time term. The
preceding tangential and wave estimates still hold. In the defect changes, the term
\NSi{NS-F-P107-024}{2V\Delta v/R} and its two contributions to the weighted flux moments are the
linear contributions to
\NSi{NS-F-P107-025}{\Delta\mathcal P,\Delta\mathcal J_\theta,\Delta\mathcal J_z} canceled by this
system. Every other term contains slow-time differentiation of the induced radial velocity, a radial
flux with \NSi{NS-F-P107-026}{b} or another radial mean, an axial derivative, radial viscosity, or an
additional tangential mean in \NSi{NS-F-P107-027}{v^2,\gamma v,\gamma^2}. These have gain greater
than \NSi{NS-F-P107-028}{0.8} above \NSi{NS-F-P107-029}{H_0}. Hence the remaining defects have
exponent greater than \NSi{NS-F-P107-030}{1.8-\kappa_s}, and both tangential means retain exponent
\NSi{NS-F-P107-031}{1.49}. These estimates imply \eqref{eq:9.8} with
\NSi{NS-F-P107-032}{B_0=0.7} and \NSi{NS-F-P107-033}{C_0^*=1.2}.

The primary curl correction has exponent \NSi{NS-F-P107-034}{1-\kappa_s>0.68}. All mean
increments and pressure changes have exponent \NSi{NS-F-P107-035}{H_0>0.9}, with radial exponent
\NSi{NS-F-P107-036}{H_0+1>1.9}. This proves \eqref{eq:9.9}. The prescribed temporal increments
have zero auxiliary mean. The first two homogeneous equations of the correction system enforce
\eqref{eq:9.10}, and the mean-field reconstruction preserves the prescribed auxiliary mean exactly.
\end{proof}

The initialization provides the bounds at stage zero. We now repeat the four corrections and compare
every new error with the bounds required at the next stage. The total corrections in \eqref{eq:9.9}
remain available when estimating interactions with the new increments.

\begin{proposition}\label{prop:9.6}
Let \NSi{NS-F-P107-037}{(u^{[j]},p^{[j]})} be a state at stage
\NSi{NS-F-P107-038}{j} as in Definition~\ref{def:9.4}. There is a state
\NSi{NS-F-P107-039}{(u^{[j+1]},p^{[j+1]})}, constructed with the same base, primary amplitudes,
supports, and inverse operators, whose residual orders are
\NSFormula{NS-F-P107-040}{display}{none}
\[
B_{j+1}=B_j+\frac1{10},
\qquad
C_{j+1}^*=C_j^*+\frac1{10}.
\]
It is obtained by the following ordered corrections:
\begin{enumerate}[label=(\roman*)]
\NSPage{108}%
\item solve the inhomogeneous amplitude equation for each supported nonzero harmonic;
\item change the wave amplitudes to correct the auxiliary-averaged tangential residual;
\item apply the temporal mean inverse to the part with zero auxiliary average;
\item apply the five-equation correction to the three current defects.
\end{enumerate}
Pressure is reconstructed after each correction. Wave and meridional velocity increments are
realized by their potentials; azimuthal mean increments are added directly. The cumulative bounds
\eqref{eq:9.9} and the exact moments \eqref{eq:9.10} are preserved.
\end{proposition}

\begin{proof}
The four steps below implement (i)--(iv) in order. For this cycle abbreviate
\NSi{NS-F-P108-001}{B=\frac12+\sigma_j} and
\NSi{NS-F-P108-002}{C_*=B+\frac12=1+\sigma_j}. At every occurrence of
\NSi{NS-F-P108-003}{g_r} and \NSi{NS-F-P108-004}{p_m} below, an asserted improved order concerns
their changes; the total corrections continue to satisfy \eqref{eq:9.9}. Within each step,
\NSi{NS-F-P108-005}{w,\beta,v,\gamma} denote the fields before that step. An increment replaces them
by
\NSFormula{NS-F-P108-006}{display}{none}
\[
w+\Delta w,
\qquad
(\beta+\Delta\beta,v+\Delta v,\gamma+\Delta\gamma).
\]
Compute \NSi{NS-F-P108-007}{g_{r,\mathrm{new}}} from these complete fields using \eqref{eq:8.3},
and set
\NSFormula{NS-F-P108-008}{display}{none}
\[
\mathcal P_{\mathrm{new}}=\int\langle g_{r,\mathrm{new}}\rangle_Y\,dR,
\qquad
p_{m,\mathrm{new}}=\mathsf T_0(g_{r,\mathrm{new}}-\rho\mathcal P_{\mathrm{new}}).
\]
All residuals and defects in the next step are computed from this updated state.

\noindent\textbf{Step 1: Cancel the supported harmonics.} We apply Proposition~\ref{prop:7.2} to
each grouped supported source, with zero initial data at the entrance to the prescribed path in the
common auxiliary torus. For source coefficient \NSi{NS-F-P108-009}{f_m}, the normalized harmonic
increments are
\NSFormula{NS-F-P108-010}{display}{none}
\[
\Delta w_m=\operatorname{curl}_*\!\left(
 \frac{i n_\Phi\times(\psi t_m)}{km|n_\Phi|^2}e^{ikm\Phi}\right),
\qquad
\Delta p_{w,m}=\psi\pi_m e^{ikm\Phi},
\qquad
\mathscr L_m(t_m,\pi_m)=-f_m.
\]
Assemble these fields with their physical rescalings over all labels and harmonics, retaining
conjugate pairs. The velocity increment has class \NSi{NS-F-P108-011}{\mathsf W_B}, and its pressure
has class \NSi{NS-F-P108-012}{\mathsf W_{B+1/2}}. The new nonzero harmonic errors have the
following exponents:
\NSFormula{NS-F-P108-013}{display}{none}
\[
\begin{array}{r|l}
\text{linear error} & B+1/2-3\kappa_s\\
\text{cross with old exact waves} & B+1/2-\kappa_s\\
\text{particular-correction self-interaction} & 2B-\kappa_s\\
\text{cross with total mean correction} & B+0.4
\end{array}
\]
The first row follows from Proposition~\ref{prop:9.1}; the other three follow from
Lemma~\ref{lem:9.2}. The mean covariance changes caused by these increments have exponent at least
\NSi{NS-F-P108-014}{B+1/2=C_*}; a radial divergence costs
\NSi{NS-F-P108-015}{\kappa_s}. The exact balances and linear pressure map therefore give
\NSFormula{NS-F-P108-016}{display}{none}
\[
E_\theta,E_z,\Delta g_r,\Delta p_m\in \mathsf M_{C_*-\kappa_s},
\qquad
(\mathcal P,\mathcal J_\theta,\mathcal J_z)\in \mathsf S_{C_*-\kappa_s}.
\]
The preserved angular-momentum and axial-flux constraints and the exact integrated identities
\eqref{eq:8.16} improve the weighted moments by one:
\NSFormula{NS-F-P108-017}{display}{9.12}
\begin{equation}
\int R^2\langle E_\theta\rangle_Y\,dR,
\qquad
\int R\langle E_z\rangle_Y\,dR
\in \mathsf S_{C_*+1-\kappa_s}.
\tag{9.12}\label{eq:9.12}
\end{equation}
The axial identity includes \NSi{NS-F-P108-018}{c_\rho\mathcal P} inside its axial derivative. This
term accounts for the normalized correction \NSi{NS-F-P108-019}{-\rho\mathcal P} left in the radial
equation after pressure reconstruction.

\noindent\textbf{Step 2: Correct the auxiliary-averaged residual by adjusting the stress.} A
compactly supported stress can cancel the auxiliary-averaged tangential residual once its total
weighted radial moments have been removed. We make this correction in physical coordinates so that
the stress agrees across charts. Use the residual components
\NSi{NS-F-P108-020}{E_\theta,E_z} of \eqref{eq:8.3}, supported in the prescribed annulus, and retain
the background residual separately with its flatness
